\section{Next Steps}
\label{sec:nextsteps}

The roadmap below is reorganized around the [D-39] empirical finding of Sec.~\ref{sec:experiments:cosmological-eval}: at the corrected anchor and the full publication schedule across four feedback physics, mean-flux and KS close while $|\Delta P_F/P_F|$ remains uniformly $\PFfactorOverGateLo\times$--$\PFfactorOverGateHi\times$ over the $\GatePFdec$ gate. The binding gate is $P_F$, the failure is structural rather than calibration, and log-space $\tau$-MSE is not a sufficient proxy for it (Sec.~\ref{sec:experiments:cosmological-eval}, ``decoupling finding''). The next-steps work is therefore on the loss / regularizer / architecture axis first, with sightline density and the 3D $\xi$ gate conditional on the structural diagnostic. The successor-work ranking in Sec.~\ref{sec:nextsteps:pf} is grounded in the [D-39] Wrinkle-1 disambiguation, now upgraded to a track-wide call by the $\tau$-binned residual decomposition of Sec.~\ref{sec:experiments:tau-binned} ([D-39] Addendum 4): the saturation band $F \in [\SatBandLo, \SatBandHi]$ absorbs a concentration ratio $R \in [\RsatRangeLo, \RsatRangeHi]$ across all four \texttt{pub-t1} cells ($> \!\RsatFactorOverFloor\times$ the positive-ID floor in $4/4$ cells, $\sim\!\RsatSpreadRelative$ relative spread cross-physics), \emph{positively identifying} saturation-regime under-fitting as the residually-supported mechanism cross-physics. The successor-work ranking below is grounded in the [D-39] Wrinkle-1 mechanism call \emph{and} the [D-40] negative result on the integrated-statistic counterfactual (Sec.~\ref{sec:nextsteps:pf} \#3): a saturation-aware $P_F$ band loss was the highest-a-priori-leverage intervention against mechanism (c) and was tested on $P_1$ at the cost-survey-equivalent schedule; eval-time $P_F$ worsened by $\PFdfortyEvalDelta$ and the flux-PDF KS gate broke. This paper's claim ceiling is therefore residually-supported mechanism, not closed-loop cause.

\begin{itemize}
    \item \textbf{Stage 1 (Foundation, done):} Sherwood ingestion, coordinate normalization, sightline + $\tau_{HI}$ loaders.
    \item \textbf{Stage 2a (Differentiable integrator, done):} RSD-convolved Voigt with learnable amplitude, re-validated at production architecture (Sec.~\ref{sec:experiments}).
    \item \textbf{Stage 2b (cost-survey + corrected-anchor):} 16-cell cost-survey micro-grid (Tab.~\ref{tab:ablation-matrix}), 12-cell Tier-1/2/3 production sweep (Sec.~\ref{sec:experiments:cost-survey-sweep}), and the four-physics \texttt{pub-t1} corrected-anchor full-schedule bundle (Tab.~\ref{tab:pub-t1-gates}). $P_F$ identified as the binding gate.
    \item \textbf{Wrinkle-1 diagnostic (priority unlock for any successor compute):} disambiguate whether the rescale-vs-trained divergence (cost-survey rescaled $P_F=\PFrescaledPone$ at $P_1$ vs.\ \texttt{pub-t1} measured $P_F=\PFresidPone$ at the same anchor) is (a) overfitting in $P_F$-relevant modes from $4\times$ more updates, (b) sightline-selection sensitivity (evaluation seed $42$ vs.\ training seed $0$), or (c) intrinsic divergence between uniform post-hoc rescale and trained-at-target in the nonlinear $F\!\leftrightarrow\!\tau$ saturation regime ([D-39]).
    \item \textbf{Stage 3 (Feedback classification, deferred):} a 4-way classifier on 3D crops of the reconstructed $\rho/\bar\rho$ field, using Stage 2b outputs as input. Out of scope for this paper --- the downstream classification claim is upstream-gated on a $P_F$-passing reconstruction, which the structural-$P_F$ gap of [D-39] places in successor work (Sec.~\ref{sec:nextsteps:stage3}).
\end{itemize}

\subsection{Closing the structural $P_F$ gap}
\label{sec:nextsteps:pf}
Three classes of intervention are warranted by the structural-$P_F$ reading. The ranking below is reordered by the [D-40] negative result on the integrated-statistic counterfactual: the previously top-ranked saturation-aware $P_F$ band loss is demoted to a negative-result entry (\#3), the per-pixel FGPA-tail regularizer is promoted to \#1, and the velocity-gradient conditioning branch holds at \#2. The W1-E trajectory (monotone $P_F$ decrease $0.5584 \to 0.4716 \to 0.4132$ at steps $5$k$/15$k$/25$k, plateau $\PFresidPone$ by step $50$k) rules out overfitting; seed sensitivity bounds sightline modulation at $\sim\!15\%$; the saturation-band concentration ratio $R \in [\RsatRangeLo, \RsatRangeHi]$ in $4/4$ cells positively identifies the saturation regime as the residually-supported mechanism cross-physics. The strict ``saturation under-fitting \emph{is the cause} of the $P_F$ gap'' claim was the target of the [D-40] counterfactual and remains open; the ranking below is actionable now.

\textbf{1.\ FGPA-tail regularizer (per-pixel physics prior on the strong-absorber regime).} The flux PDF closes in $3/4$ cells; the $P_2$ KS breach ($\KSseedFortyTwoPtwo$) and the uniform $P_F$ gap suggest the tail behavior of $F$ (the strong-absorber regime, $F \to 0$) is the unconstrained mode --- the same band the $\tau$-binned decomposition of Sec.~\ref{sec:experiments:tau-binned} localizes the residual mass to ($R \geq \RsatRangeLo$ in all four cells). A regularizer that penalizes deviation of $\hat F$'s tail moments from the FGPA prediction $\tau \propto \rho^{1.6}\,T^{-0.7}$~\cite{bi_davidsen1997,hui_gnedin1997} constrains the $\tau$-PDF tails where the saturation-regime under-fitting positively identified track-wide is concentrated. Unlike the integrated-statistic loss of \#3, this constraint is enforced \emph{per pixel} along the rendered sightline --- the FGPA scaling fixes a local $\tau(\rho, T)$ relation that the network must respect at every sample, not merely in the band-integrated mean. This is the structural property that makes \#1 immune to the constant-prediction degeneracy that broke \#3: a uniform $P_{F,\text{pred}}\!\to\!0$ collapse cannot simultaneously satisfy a per-pixel $\tau\!\propto\!\rho^{1.6}T^{-0.7}$ prior driven by the predicted $(\rho, T)$ fields. The regularizer is differentiable, requires no new inputs, and is proposed as a design entry --- empirical leverage on the [D-39] mechanism is not claimed in advance of a run.

\textbf{2.\ Velocity-gradient conditioning branch.} A conditioning branch that feeds finite-difference estimates of $\partial v_{\text{pec}}/\partial \chi$ along the line of sight into the density head is proposed as the next test of the ground-truth-anchored anti-degeneracy discipline ([D-42-meta], [D-37] extension). The motivation is that the Voigt kernel has explicit $v_{\text{pec}}$ dependence but $F_\theta$ is currently unconditional on local velocity-gradient structure, and that the conditioning input is observation-derived rather than network-predicted --- the cleanest known way to avoid the self-consistent-state degeneracy that broke the integrated-statistic loss of \#3 and the per-pixel FGPA-tail prior of \#1 by allowing the network to game the regularizer through constant-prediction collapse. This is not a high-confidence pick: two prior \#1/\#2 specs at similar a-priori confidence have been empirically falsified by previously-unaudited degeneracies, and the present design has only passed the cleanest pre-run check we currently know, not an actual run. It is not claimed to address a different failure mode than the loss-side interventions; it may address the same saturation-band under-fitting from a different axis (RSD line-shape conditioning rather than tail-moment penalty), which the actual run will test.

\textbf{3.\ Saturation-aware $P_F$ band loss --- negative result ([D-40]).} A saturation-aware $P_F$ band loss was tested as the highest-a-priori-leverage intervention against the [D-39] mechanism-(c) finding. On $P_1$ at the cost-survey-equivalent schedule ($12{,}500$ steps, $\langle F\rangle_{\text{obs}}=\AnchorMeanF$, $n_{\text{rays}}=64$, seed $=0$, weights \texttt{sat\_band\_weight}$=3.0$, \texttt{rank\_order\_weight}$=0.1$, \texttt{pf\_loss\_weight}$=1.0$), the band-integrated relative-residual term descended from $0.99$ at initialization to $\sim\!7\times\!10^{-6}$ at step $10{,}000$ in training. Eval-time measurement against the same fiducial point produced $|\Delta P_F/P_F|=\PFdfortyEval$ ($+\PFdfortyEvalDelta$ worse than the [D-39] \texttt{pub-t1} $P_1$ baseline of $\PFresidPone$) and broke the flux-PDF KS gate ($\KSdfortyEval$ vs the $\GateKS$ bar); training $\log(1+\tau)$-MSE worsened $4\times$ relative to the cost-survey baseline at the same step count. The integrated relative-residual formulation $((P_{F,\text{pred}}-P_{F,\text{truth}})/P_{F,\text{truth}})^2$ admits trivial constant-prediction solutions: when $P_{F,\text{truth}}$ is small in any band, the relative residual is minimized by driving $P_{F,\text{pred}}\!\to\!0$ uniformly at the cost of all data fit. This motivates the per-pixel FGPA-tail constraint of \#1 above, which is structurally immune to the integrated-statistic degeneracy. See LEDGER [D-40].

\subsection{Tier 4 and sightline-density scaling}
\label{sec:nextsteps:tier4}
Tier 4 (sightline density $n_{\text{rays}}=16{,}384$, $\sim\!\GPUHoursTierFour$ GPU-hours across four feedback variants) is the natural completion of the degradation matrix of Tab.~\ref{tab:ablation-matrix} and the unlock for the full 3D $\xi_{\hat\rho,\rho}(r)$ gate (the \texttt{SherwoodIGM\_gal} dataset has landed on the Juno scratch tier and the evaluator is ready). Tier 4 dispatch is \emph{conditional on the Wrinkle-1 mechanism outcome}, not a forward commitment to close $P_F$. With the $\tau$-binned decomposition (Sec.~\ref{sec:experiments:tau-binned}) positively identifying saturation-regime under-fitting cross-physics, the $\sim\!\GPUHoursTierFour$-GPU-hr budget would not be expected to close $P_F$ on sightline density alone --- a saturation deficit that is physics-invariant across $4/4$ cells at fixed $n_{\text{rays}}=64$ is not a function the dense-sightline axis is positioned to address. The priority therefore returns to the loss / regularizer / architecture axis of Sec.~\ref{sec:nextsteps:pf}, and Tier 4 is sequenced behind the promoted FGPA-tail regularizer (\#1) following the [D-40] negative result on the integrated-statistic counterfactual.

The argument-from-invariance bears explicit unpacking, because the deferral rests on it. The four feedback variants of the \texttt{pub-t1} bundle --- $P_1$ (no feedback), $P_2$ (winds), $P_3$ (winds$+$AGN), $P_4$ (strong AGN) --- span the full range of small-scale gas thermodynamics available in the Sherwood suite, yet at fixed $n_{\text{rays}}=64$ all four cells fail $P_F$ at the same $\PFfactorOverGateLo\times$--$\PFfactorOverGateHi\times$ over the $\GatePFdec$ gate, with the saturation-band concentration ratio $R \in [\RsatRangeLo, \RsatRangeHi]$ in $4/4$ cells ([D-39], cross-physics bootstrap-confirmed under [D-44]). A deficit that is invariant under the largest available physics perturbation is unlikely to be a function the sightline-density axis is positioned to address, for two reasons. (i) Sightline density modulates the \emph{sampling} of an underlying field, not the \emph{form} of that field --- the saturation-regime under-fitting positively identified track-wide is a property of how $F_\theta$ represents the $F\!\in\![\SatBandLo, \SatBandHi]$ band, not of how many sightlines that band is observed through. (ii) If the model does not recover the saturated absorber regime from $64$ rays under any of the four physics, adding samples of the same physics would not be expected to change the loss landscape's failure mode --- more samples constrain the parameters of the existing representation, they do not alter the representation's expressive capacity in the saturation regime. Tier 4 is therefore not the cheapest \emph{informative} next experiment on the present evidence; the three retired counterfactual interventions ([D-40] integrated band loss, [D-41] FGPA-tail per-pixel regularizer, [D-42] velocity-gradient conditioning) point uniformly to the loss / regularizer / architecture axis as the locus where the binding $P_F$ deficit is expected to move. We retain the explicit deferral over a sightline-density commitment; if a reviewer demands empirical Tier-4 evidence, a single-cell eval-only pilot at $n_{\text{rays}}=16{,}384$ on one physics ($\sim\!\$1.50$ Juno spend) would settle the question without changing the paper's verdict in either direction.

\subsection{Stage 3: Feedback Identification from Reconstructed $\rho$}
\label{sec:nextsteps:stage3}
Stage 3 trains a four-way classifier on 3D crops of the reconstructed $\rho/\bar\rho$ field over $\{\text{P1}, \text{P2}, \text{P3}, \text{P4}\}$, with target accuracy $> 85\%$. The data-pipeline contract is fixed: input is a $(C{=}1, D, H, W)$ cube of reconstructed overdensity at native simulation resolution, label is the originating \texttt{physics\_id}. The architectural choice (3D CNN vs.\ 3D ViT, crop dimensions, augmentation policy) is deferred to a Stage 3 design pass and is gated on a $P_F$-passing reconstruction; the structural-$P_F$ gap of [D-39] is upstream of the feedback-discrimination claim.

\subsection{Substrate-scale follow-on: first realized 48$^3$ probe}
\label{sec:nextsteps:substrate-cprime}
The substrate-discriminability probe at $\SubstrateProbeCropSize^3$ truth-$\rho$ crops (suppl.) declined to establish that the substrate at that configuration carries cross-physics signal beyond two summary moments, with a 3D ResNet--mean+variance gap of $\SubstrateProbeMarginPP$~pp against a $\SubstrateProbeADfiveBarPP$~pp self-anchored reference. The first realized substrate-scale follow-on along that axis re-runs the probe on $\SubstrateProbeCpCropSize^3$ truth-$\rho$ crops at \emph{fixed} simulation pitch ($n_{\text{grid}}=768$, per-voxel comoving edge $\SubstrateProbeCpVoxelKpc$~kpc/h unchanged); only the crop dimension changes, so the comparison is a clean substrate-extent contrast rather than a resolution contrast~\cite{touvron2019fixres}. The 4-class Sherwood feedback discrimination signature at $\SubstrateProbeCpCropSize^3$ exceeds the moment-only MVSK 4-scalar subspace by $\SubstrateProbeCpMarginPP$~pp ($5$-seed mean; $95\%$ bootstrap CI $[\SubstrateProbeCpCIlow, \SubstrateProbeCpCIhigh]$~pp; $5/5$ seeds clear the self-anchored $\SubstrateProbeCpRefPP$~pp reference). The reading is therefore complementary to but not a generalization of the $\SubstrateProbeCropSize^3$ result: at $\SubstrateProbeCropSize^3$ the same problem collapses to the MVSK 4-scalar subspace within seed noise; at $\SubstrateProbeCpCropSize^3$ it does not. This $(\SubstrateProbeCropSize^3, \SubstrateProbeCpCropSize^3)$ pair, at fixed pitch, establishes a single-pair substrate-scale contrast~\cite{alain_bengio2017probing} --- substrate scale matters in this direction at this configuration. No curve-shape claim is admissible from two points; a multi-point scoping curve ($96^3$+) is the natural next step.

The empirical seed standard deviation on the 3D ResNet readout, $\SubstrateProbeCpResnetSigma$, exceeded the pre-committed seed-variance budget ($\SubstrateProbeCpSeedBudget$) by a factor of $\SubstrateProbeCpSeedFactor$; the effect size at $\SubstrateProbeCpCropSize^3$ was large enough that this under-estimate did not flip the verdict, but it confirms the directional risk on discriminative-classifier seed sensitivity~\cite{damour2020underspecification} in the conservative direction. The MVSK baseline that the ResNet is compared against here is nested inside the same self-anchored reference framework as the $\SubstrateProbeCropSize^3$ probe (suppl.), and the $\SubstrateProbeCpRefPP$~pp reference is project-side --- it is not a published external threshold. Both the project-side bar and the nesting are disclosed so that the substrate-axis reading is not overdrawn.

\subsection{Method decision log}
The methodology of Sec.~\ref{sec:method} and the evaluation contract of Sec.~\ref{sec:experiments} are the consolidated form of a longer decision sequence. Tab.~\ref{tab:decision-log} summarizes the scientific entries; full traceability lives in the project's experiment LEDGER.

\begin{table}[t]
    \centering
    \renewcommand{\arraystretch}{1.15}
    \setlength{\tabcolsep}{4pt}
    \begin{tabular}{l|p{0.74\columnwidth}}
    ID & Decision \\
    \hline
    {[D-01]} & Tepper-Garc\'ia 4th-order Voigt; small-$|x|$ Taylor branch added 2026-05-04 \\
    {[D-02]} & Bounded-activation physics layers for $\rho$, $T$, $X_{HI}$, $v_{\text{pec}}$ \\
    {[D-06]} & RSD-convolved integrator (Stage 2a$\to$2b lift); first-half = redshift-space target \\
    {[D-11]/[D-34]} & Mean-flux soft anchor $\langle F\rangle_{\text{obs}} = \AnchorMeanF \pm \AnchorMeanFsigma$ at $z=0.3$ (Kirkman+ 2007); cost-survey runs trained against retracted value $0.877$, kept under anchor-invariance \\
    {[D-13]} & Stage 2b gates: $P_F(k_\parallel)$, Pearson $\xi_{\hat\rho,\rho}(2\,h^{-1}\,\text{Mpc})$, flux-PDF KS on $F\in[\PDFwinLo,\PDFwinHi]$ \\
    {[D-14]} & Production schedule (per-tier max\_steps, microbatch, accum) \\
    {[D-19]} & Smoke-pass criteria: gradient finite, descent $\leq 0.85$, mean\_F within $5\%$ \\
    {[D-21]} & Two-pass surrogate gradient for the masked mean-flux constraint \\
    {[D-23]} & Tier-aware microbatch / step schedule with linear VRAM model \\
    {[D-24]} & log1p+cap+mask Ly$\alpha$ forest loss; $\tau_{\max}=10$ LOCKED 2026-05-04 \\
    {[D-35]} & Cost-survey post-hoc rescale read as anchor-preview, retained under [D-39] only as upper-bound optimism \\
    {[D-38]} & log-space supervision is the methods contribution; cap and mask retained as defense-in-depth \\
    {[D-39]} & \texttt{pub-t1} corrected-anchor full-schedule bundle falsifies the rescale-preview; $P_F$ is the structurally binding gate; Tier 4 conditional on Wrinkle-1 \\
    {[D-40]} & Saturation-aware $P_F$ band loss tested on $P_1$ (12.5k-step counterfactual): eval $|\Delta P_F/P_F|=\PFdfortyEval$ ($+37\%$ vs.\ baseline), KS $=\KSdfortyEval$ (gate broken); integrated-residual formulation admits constant-prediction degeneracy. \#4.1 ranking re-pivoted to per-pixel FGPA-tail prior. \\
    \end{tabular}
    \caption{Methodological decision log. One-line summary per scientific entry; LEDGER \S3 carries full derivations, sources, and verification status. [D-39] supersedes the verdict-frame implicit in earlier decision entries: the corrected-anchor full-schedule run is recorded as a calibration-vs-structure falsification, not a thresholded PASS.}
    \label{tab:decision-log}
\end{table}
